%% kapa: exam=1 round=22 year=2011 period=1 form=A subject=economics subject_ko=경제원론 seq=36 number=76
%% source: Q-Net 감정평가사 기출문제 (cst00309) artlSeq=5036790 file=extracted/1008/5036790/2011년 제22회 감정평가사 1차시험 기출문제/1교시.hwp
%% conversion: kordoc 4.24.51 -> tools/gen1.py (auto); review=proofread-sonnet; flags=none
\begin{question}{76}
외부효과에 관한 설명으로 옳은 것은?
\begin{choices}
  \item 생산의 외부불경제가 존재하는 경우 사회적 최적생산량은 시장균형생산량보다 많다.
  \item 소비의 외부경제가 존재하는 경우 사회적 최적소비량은 시장균형소비량보다 적다.
  \item 외부효과의 내부화로는 외부효과의 비효율성을 해결할 수 없다.
  \item 교정적 조세는 경제적 효율을 향상시키면서 정부의 조세수입도 증대시킨다.
  \item 오염배출권 거래제에서는 정부가 오염배출권의 가격을 먼저 설정함으로써 사회적 총오염배출량이 결정된다.
\end{choices}
\end{question}
\answer{4}
